\documentclass[12pt, oneside]{book}
\usepackage{geometry}
\geometry{margin=1.25in}

% ─── Fonts ───────────────────────────────
\usepackage{fontspec}
\setmainfont{Padauk}

% ─── Math ────────────────────────────────
\usepackage{amsmath, amssymb, bm}

% ─── Tables ──────────────────────────────
\usepackage{booktabs, array}

% ─── Misc ────────────────────────────────
\usepackage{graphicx, microtype}
\usepackage{hyperref}

\usepackage{tcolorbox}
\usepackage{fontawesome5}

\definecolor{infocolor}{RGB}{219, 234, 254}
\definecolor{infoborder}{RGB}{59, 130, 246}

\newtcolorbox{infobox}{
  colback=infocolor,
  colframe=infoborder,
  arc=4pt,
  boxrule=1pt,
  left=6pt,
  right=6pt,
  top=4pt,
  bottom=4pt,
  title={{\faYoutube \quad သိထားသင့်တာလေးများ}},
  coltitle=white
}

\title{Neural Networks with Maths\\[0.3em]
       \large သင်္ချာအခြေခံများဖြင့်}
\author{Light}
\date{}

\begin{document}
\maketitle

\chapter{သင်္ချာအခြေခံများ}

Neural network ဆိုတာ သင်္ချာညီမျှခြင်းများတွေ ပြေလည်စေမယ့် အဖြေများရှာခြင်းဖြစ်တယ်။ \\
အောက်က ဥပမာတွေကို ကြည့်ပါ။ ပေးထားချက် (Input) နှစ်ခု $x_1$, $x_2$ နှင့် အဖြေ (output) $y$ ရှိတယ်။ အဲတာကို ဥပမာ ၃ကြောင်းပေးထားတယ်။

\begin{center}
\begin{tabular}{ccc}
\toprule
$x_1$ & $x_2$ & $y$ \\
\midrule
1 & 2 & 5 \\
1 & 3 & 7 \\
2 & 1 & 4 \\
\bottomrule
\end{tabular}
\end{center}

$$
w_1 x_1 + w_2 x_2 = y
$$
ဒီပြဿနာရဲ့ ညီမျှခြင်းက ဒီလိုဆိုပါတော့။ အဲတာဆိုရင် အပေါ်က ဥပမာ ၃ခုလုံးကို မှန်စေမဲ့ $w_1$, $w_2$ တန်ဖိုးတွေကို ရှာကြမယ်။

\begin{align}
w_1 (1) + w_2 (2) = 5 \\
w_1 (1) + w_2 (3) = 7 \\
w_1 (2) + w_2 (1) = 4
\end{align}

\noindent ဒါဆိုရင် အခု ညီမျှခြင်း ၃ကြောင်းလုံးကို မှန်စေမဲ့ $w_1$, $w_2$ တန်ဖိုးတွေကို တွက်ကြည့်။

\noindent ညီမျှခြင်း $(1.2)$ နှင့် $(1.1)$ ကိုနှုတ်ကြည့်ရင် $w_2 = 2$ ရမယ်။ အဲတာဆိုရင် $$ w_1 (1) + 2 (2) = 5 \implies w_1 = 1 $$

\section{Linear Algebra အခြေခံ}
အခု ပုစ္ဆာကိုပဲ matrix တွေနဲ့ တွက်ကြည့်ရအောင်။ ညီမျှခြင်း $(1.1)$ နှင့် $(1.2)$ နှစ်ကြောင်းထဲကို ယူပြီး matrix form နဲ့ ရေးမယ်။
$$
\begin{bmatrix} 1 & 2 \\ 1 & 3 \end{bmatrix} \begin{bmatrix} w_1 \\ w_2 \end{bmatrix} = \begin{bmatrix} 5 \\ 7 \end{bmatrix}
$$
ဒါကို $X\mathbf{w} = \mathbf{y}$ လို့ ရေးနိုင်တယ်။ ဒီမှာ
$$
X = \begin{bmatrix} 1 & 2 \\ 1 & 3 \end{bmatrix}, \quad 
\mathbf{w} = \begin{bmatrix} w_1 \\ w_2 \end{bmatrix}, \quad 
\mathbf{y} = \begin{bmatrix} 5 \\ 7 \end{bmatrix}
$$
\noindent $\mathbf{w}$ ကို ရှာဖို့ နှစ်ဖက်လုံးကို $X^{-1}$ နဲ့ မြှောက်မယ်။
$$
X^{-1} X \mathbf{w} = X^{-1} \mathbf{y}
$$
\noindent $X^{-1}X = I$ ဖြစ်တယ်။ $I$ ကို Identity Matrix လို့ခေါ်တယ်။
$$
I = \begin{bmatrix} 1 & 0 \\ 0 & 1 \end{bmatrix}
$$
\noindent ဘယ် vector ကိုမဆို $I$ နဲ့ မြှောက်ရင် မပြောင်းဘဲ ဒီအတိုင်းနေတယ်။ ဆိုလိုသည်မှာ $I\mathbf{w} = \mathbf{w}$ ဖြစ်တယ်။ လွယ်လွယ်ပြောရင် ကျေသွားတယ်ပေါ့။ ဒါကြောင့်
$$
\mathbf{w} = X^{-1} \mathbf{y}
$$
\noindent ခုဆိုရင် $X^{-1}$ ကို ရှာဖို့သာ လိုတော့တယ်။ $2 \times 2$ matrix အတွက် inverse formula က
$$
\begin{bmatrix} a & b \\ c & d \end{bmatrix}^{-1} = \frac{1}{ad - bc} \begin{bmatrix} d & -b \\ -c & a \end{bmatrix}
$$
\begin{align*}
X^{-1} &= \frac{1}{(1)(3) - (2)(1)} \begin{bmatrix} 3 & -2 \\ -1 & 1 \end{bmatrix} = \begin{bmatrix} 3 & -2 \\ -1 & 1 \end{bmatrix} \\[10pt]
\mathbf{w} &= X^{-1}\mathbf{y} = \begin{bmatrix} 3 & -2 \\ -1 & 1 \end{bmatrix} \begin{bmatrix} 5 \\ 7 \end{bmatrix} = \begin{bmatrix} (3)(5) + (-2)(7) \\ (-1)(5) + (1)(7) \end{bmatrix} = \begin{bmatrix} 1 \\ 2 \end{bmatrix}
\end{align*}
\noindent အဖြေက $w_1 = 1$, $w_2 = 2$

\begin{infobox}
Linear Algebra ကို အခြေခံပြန်ကြည့်ချင်ရင် 3Blue1Brown ရဲ့ series ကို သွားကြည့်ဖို့ အကြံပေးပါတယ်။ \\
\href{https://youtube.com/playlist?list=PLZHQObOWTQDPD3MizzM2xVFitgF8hE_ab&si=QxWZPVdOINlLmugE}{3Blue1Brown --- Essence of Linear Algebra}
\end{infobox}
\end{document}
